\section{Introduction}

The direct formal-exponent marker of \cite{RuddM1} proves that
\(X\times\PP^1\) is irrational for every smooth cubic threefold \(X\).
This note studies a finer count on the same generic even quantum
\(D\)-module.  It retains the marked small section and the original loop
framing and counts primitive-sixth formal-monodromy eigenvalues.  We denote
the resulting count by \(\nu_6\).

The cubic calculation and the product formula are unconditional:
\[
 \nu_6(X)=2,\qquad \nu_6(X\times\PP^1)=4.
\]
The operation formulas require one additional reconstruction statement, and
the remaining low-dimensional center argument requires a separate
divisor-tagging statement for surface centers that are neither minimal nor
geometrically ruled.  We state both precisely where they enter.

\begin{theorem}[Conditional framed obstruction]
\label{thm:every-cubic-conditional}
\coverage{conditional_deduction}%
\lean{cubicThreefold_oneProjectiveLine_conclusion_of_framedMarker}%
Assume Hypothesis~R below, and Hypothesis~T only for surface centers
that are neither minimal nor geometrically ruled.  For every smooth complex
cubic threefold \(X\),
\[
 \nu_6(X\times\PP^1)=4,
\]
every rational smooth projective fourfold has \(\nu_6=0\), and consequently
\(X\times\PP^1\) is irrational.
\end{theorem}

The conclusion is not needed for the unconditional theorem of
\cite{RuddM1}.  Its purpose is to isolate a finer invariant and the exact
geometric comparison inputs still needed to make its birational invariance
unconditional.

Section~\ref{sec:framed-monodromy} defines the count, proves the
unconditional cubic and product calculations, states the two residual
hypotheses, proves conditional operation formulas and low-dimensional
vanishing, and derives the theorem.
